\usepackage{amsmath,amssymb}             % AMS Math
\usepackage{amsthm}                      % Entornos de teorema y demostración
% \usepackage[french]{babel}

% Entornos de teorema, numerados de forma conjunta por capítulo
\theoremstyle{plain}
\newtheorem{teorema}{Teorema}[chapter]
\newtheorem{proposicion}[teorema]{Proposición}
\renewcommand{\proofname}{Demostración}
\usepackage[T1]{fontenc}
\usepackage[left=1.0in,right=1.0in,top=1.0in,bottom=1.0in,includefoot,includehead,headheight=14pt]{geometry}
\renewcommand{\baselinestretch}{1.05}
%comando faltante
\usepackage{booktabs,array,graphicx,needspace}
\usepackage{tikz}
\usetikzlibrary{arrows.meta,positioning}
\setlength{\emergencystretch}{3em}
% Table of contents for each chapter

\usepackage[nottoc]{tocbibind}           % lista de figuras y de tablas, en el índice

% babel en español rotula los flotantes «table» como «Cuadro»; este documento usa «Tabla»
\addto\captionsspanish{%
  \renewcommand{\tablename}{Tabla}%
  \renewcommand{\listtablename}{Índice de tablas}%
}
\usepackage{minitoc}
\setcounter{minitocdepth}{2}
\mtcindent=15pt
% Use \minitoc where to put a table of contents

\usepackage{aecompl}

% My pdf code

\usepackage{ifpdf}


% Links in pdf
\usepackage{color}
\definecolor{linkcol}{rgb}{0,0,0.4} 
\definecolor{citecol}{rgb}{0.5,0,0} 

% Change this to change the informations included in the pdf file

% See hyperref documentation for information on those parameters



% definitions.
% -------------------

\setcounter{secnumdepth}{3}
\setcounter{tocdepth}{2}

% Some useful commands and shortcut for maths:  partial derivative and stuff

\newcommand{\pd}[2]{\frac{\partial #1}{\partial #2}}
\def\abs{\operatorname{abs}}
\def\argmax{\operatornamewithlimits{arg\,max}}
\def\argmin{\operatornamewithlimits{arg\,min}}
\def\diag{\operatorname{Diag}}
\newcommand{\eqRef}[1]{(\ref{#1})}

\usepackage{rotating}                    % Sideways of figures & tables
%\usepackage{bibunits}
%\usepackage[sectionbib]{chapterbib}          % Cross-reference package (Natural BiB)
%\usepackage{natbib}                  % Put References at the end of each chapter
                                         % Do not put 'sectionbib' option here.
                                         % Sectionbib option in 'natbib' will do.
\usepackage{fancyhdr}                    % Fancy Header and Footer

% \usepackage{txfonts}                     % Public Times New Roman text & math font
  
%%% Fancy Header %%%%%%%%%%%%%%%%%%%%%%%%%%%%%%%%%%%%%%%%%%%%%%%%%%%%%%%%%%%%%%%%%%
% Fancy Header Style Options

\pagestyle{fancy}                       % Sets fancy header and footer
\fancyfoot{}                            % Delete current footer settings

%\renewcommand{\chaptermark}[1]{         % Lower Case Chapter marker style
%  \markboth{\chaptername\ \thechapter.\ #1}}{}} %

%\renewcommand{\sectionmark}[1]{         % Lower case Section marker style
%  \markright{\thesection.\ #1}}         %

\fancyhead[LE,RO]{\bfseries\thepage}    % Page number (boldface) in left on even
% pages and right on odd pages
\fancyhead[RE]{\bfseries\nouppercase{\leftmark}}      % Chapter in the right on even pages
\fancyhead[LO]{\bfseries\nouppercase{\rightmark}}     % Section in the left on odd pages

\let\headruleORIG\headrule
\renewcommand{\headrule}{\color{black} \headruleORIG}
\renewcommand{\headrulewidth}{1.0pt}
\usepackage{colortbl}
\arrayrulecolor{black}

\fancypagestyle{plain}{
  \fancyhead{}
  \fancyfoot{}
  \renewcommand{\headrulewidth}{0pt}
}

\usepackage{algorithm}
\usepackage[noend]{algorithmic}

% Pseudocódigo en español
\floatname{algorithm}{Algoritmo}
\renewcommand{\listalgorithmname}{Índice de algoritmos}
\renewcommand{\algorithmicrequire}{\textbf{Entrada:}}
\renewcommand{\algorithmicensure}{\textbf{Salida:}}
\renewcommand{\algorithmicfor}{\textbf{para}}
\renewcommand{\algorithmicforall}{\textbf{para todo}}
\renewcommand{\algorithmicdo}{\textbf{hacer}}
\renewcommand{\algorithmicwhile}{\textbf{mientras}}
\renewcommand{\algorithmicif}{\textbf{si}}
\renewcommand{\algorithmicthen}{\textbf{entonces}}
\renewcommand{\algorithmicelse}{\textbf{si no}}
\renewcommand{\algorithmicend}{\textbf{fin}}
\renewcommand{\algorithmicreturn}{\textbf{devolver}}

%%% Clear Header %%%%%%%%%%%%%%%%%%%%%%%%%%%%%%%%%%%%%%%%%%%%%%%%%%%%%%%%%%%%%%%%%%
% Clear Header Style on the Last Empty Odd pages
\makeatletter

\def\cleardoublepage{\clearpage\if@twoside \ifodd\c@page\else%
  \hbox{}%
  \thispagestyle{empty}%              % Empty header styles
  \newpage%
  \if@twocolumn\hbox{}\newpage\fi\fi\fi}

\makeatother
 
%%%%%%%%%%%%%%%%%%%%%%%%%%%%%%%%%%%%%%%%%%%%%%%%%%%%%%%%%%%%%%%%%%%%%%%%%%%%%%% 
% Prints your review date and 'Draft Version' (From Josullvn, CS, CMU)
\newcommand{\reviewtimetoday}[2]{\special{!userdict begin
    /bop-hook{gsave 20 710 translate 45 rotate 0.8 setgray
      /Times-Roman findfont 12 scalefont setfont 0 0   moveto (#1) show
      0 -12 moveto (#2) show grestore}def end}}
% You can turn on or off this option.
% \reviewtimetoday{\today}{Draft Version}
%%%%%%%%%%%%%%%%%%%%%%%%%%%%%%%%%%%%%%%%%%%%%%%%%%%%%%%%%%%%%%%%%%%%%%%%%%%%%%% 

\newenvironment{maxime}[1]
{
\vspace*{0cm}
\hfill
\begin{minipage}{0.5\textwidth}%
%\rule[0.5ex]{\textwidth}{0.1mm}\\%
\hrulefill $\:$ {\bf #1}\\
%\vspace*{-0.25cm}
\it 
}%
{%

\hrulefill
\vspace*{0.5cm}%
\end{minipage}
}

\let\minitocORIG\minitoc
\renewcommand{\minitoc}{\minitocORIG \vspace{1.5em}}

\usepackage{multirow}
\usepackage{slashbox}

\newenvironment{bulletList}%
{ \begin{list}%
	{$\bullet$}%
	{\setlength{\labelwidth}{25pt}%
	 \setlength{\leftmargin}{30pt}%
	 \setlength{\itemsep}{\parsep}}}%
{ \end{list} }

\renewcommand{\epsilon}{\varepsilon}

% centered page environment

\newenvironment{vcenterpage}
{\newpage\vspace*{\fill}\fancyhf{}\renewcommand{\headrulewidth}{0pt}}
{\vspace*{\fill}\par\pagebreak}
% Cargar después de los paquetes de flotantes.
\usepackage{hyperref}
\hypersetup
{
bookmarksopen=true,
pdftitle={Predicción temprana del abandono en LMS: versión de trabajo},
pdfauthor={Oriana Giraldo Arcia y Luis Javier Rubio Hernández}, 
pdfsubject={Proyecto Aplicado de Ciencia de Datos}, %subject of the document
%pdftoolbar=false, % toolbar hidden
pdfmenubar=true, %menubar shown
pdfhighlight=/O, %effect of clicking on a link
colorlinks=true, %couleurs sur les liens hypertextes
pdfpagemode=UseNone, %aucun mode de page
pdfpagelayout=SinglePage, %ouverture en simple page
pdffitwindow=true, %pages ouvertes entierement dans toute la fenetre
linkcolor=linkcol, %couleur des liens hypertextes internes
citecolor=citecol, %couleur des liens pour les citations
urlcolor=linkcol %couleur des liens pour les url
}

% Estilo visual del proyecto, basado en la identidad institucional.
% Generado desde estilo_visual_javeriana.json; no cambia el cuerpo ni las ecuaciones.
\definecolor{PUJAzul}{HTML}{2C5697}
\definecolor{PujAcento}{HTML}{347B83}
\definecolor{PUJGris}{HTML}{B8BDC5}
\definecolor{PUJTexto}{HTML}{202020}
\colorlet{PUJClaro}{PUJAzul!6!white}
\colorlet{PUJResalte}{PujAcento!18!white}
% etoolbox ya está disponible mediante hyperref; sin nuevas dependencias tipográficas.
\makeatletter
\patchcmd{\@makechapterhead}{\normalfont}{\normalfont\color{PUJAzul}}{}{}
\patchcmd{\@makeschapterhead}{\normalfont}{\normalfont\color{PUJAzul}}{}{}
\patchcmd{\section}{\normalfont\Large\bfseries}{\normalfont\Large\bfseries\color{PUJAzul}}{}{}
\patchcmd{\subsection}{\normalfont\large\bfseries}{\normalfont\large\bfseries\color{PUJAzul}}{}{}
\patchcmd{\subsubsection}{\normalfont\normalsize\bfseries}{\normalfont\normalsize\bfseries\color{PUJAzul}}{}{}
\makeatother
\colorlet{linkcol}{PUJAzul}
\colorlet{citecol}{PUJAzul}
\arrayrulecolor{PUJAzul}
\renewcommand{\headrule}{\color{PUJAzul}\headruleORIG}
\tikzset{puj base/.style={draw=PUJAzul,text=PUJTexto},
  puj etapa/.style={draw=PUJAzul,text=PUJTexto,rounded corners,line width=.6pt}}

